% ============================================================
% Recreational Marijuana Legalization and Drug-Overdose
% Mortality: A Cross-Border Distance Design
% Formatted for ECONOMICS LETTERS (Elsevier, elsarticle class).
% REVISION (2026-07-26): adds a dispensary-level robustness check
% that replaces the main design's distance-to-nearest-legal-state
% -county measure with distance to the nearest ACTUALLY OPERATING
% dispensary, using real license/opening dates geocoded from state
% cannabis-regulator records (9 of 24 legal states covered). The
% original state-legal-county-centroid design, and every earlier
% check (pre-trends, medical placebo, persistence, county FE,
% period-specific effects, welfare), is retained as the paper's
% primary specification; the new check is reported as an additional
% robustness section, continuous distance kept as the primary
% quantitative measure throughout, consistent with the rest of the
% paper.
% Compile: pdfLaTeX -> BibTeX -> pdfLaTeX -> pdfLaTeX
% ============================================================
\documentclass[preprint,12pt,authoryear]{elsarticle}

\usepackage{amsmath,amssymb}
\usepackage{graphicx}
\usepackage{booktabs}
\usepackage{enumitem}
\usepackage[hidelinks]{hyperref}

\journal{Economics Letters}

\begin{document}

\begin{frontmatter}

\title{Recreational Marijuana Legalization and Drug-Overdose Mortality:\\
A Cross-Border Distance Design}

\author{Author name(s) withheld for review}

\begin{abstract}
\noindent We test whether proximity to a legal recreational-marijuana
market reduces drug-overdose mortality, using a stacked multi-period
panel of counties in 31 never-legalizing U.S. states and every state's
actual legalization date, 2020--2025, rather than a single
before/after snapshot. Switcher counties (distance to the nearest
legal-market county crossing below 50 miles) see overdose deaths fall
2.32 more per 100,000 than controls in the same transition
($p=0.004$); a continuous per-mile specification agrees in sign and is
itself significant ($p<0.001$). A formal pre-trends test shows the raw
switcher comparison is contaminated by mean reversion, which a
baseline-rate control corrects; the continuous specification, unlike
the switcher indicator, remains significant under state and county
fixed effects, though both specifications lose significance under the
sparser state-by-transition interaction, which we read as
low-power/overfit rather than a genuine reversal. A placebo check using
medical- rather than recreational-marijuana timing is null. The effect
does not compound in years after switching. Comparing our stacked
design against a canonical two-way fixed-effects regression and the
\citet{CallawaySantAnna2021} group-time estimator on the same
treatment definition, the three do not closely agree: two-way fixed
effects roughly doubles our estimate, in the direction predicted by
its known bias under staggered adoption, while Callaway--Sant'Anna
returns a smaller, insignificant estimate whose own event-study
reveals the identical pre-treatment mean-reversion problem we
independently diagnose and correct for. As a further robustness
check, we replace the state-legal-county distance measure with
distance to the nearest \emph{actually operating} dispensary, using
real license-issuance dates geocoded from state cannabis-regulator
records for 9 of the 24 legal states; the continuous coefficient stays
positive and of comparable magnitude, though only marginally
significant given the smaller, partial-coverage sample. Using
already-near counties as an alternative control group instead of
far/never-treated counties, on both treatment-timing definitions: with
state-legal-county distance, the switcher coefficient is never
significant, even without controls; with distance to a real, operating
dispensary, it is significant without controls but falls by more than
half and loses significance once the same mean-reversion correction,
or state or county fixed effects, are applied. Applying the national
age distribution of
overdose deaths, the switcher effect implies roughly 755 averted deaths
and 22,500 years of life saved -- illustrative, not confirmed.
\end{abstract}

\begin{keyword}
recreational marijuana \sep opioid overdose mortality \sep cross-border spillover \sep staggered difference-in-differences \sep JEL: I18, I12, K42
\end{keyword}

\end{frontmatter}

\vspace{0.3cm}
\noindent\textbf{Highlights}
\begin{itemize}[leftmargin=1.4em]
\item Stacked multi-period cross-border distance design, using every state's actual legalization date, 2020--2025.
\item Switcher counties show 2.32 fewer deaths/100k ($p{=}0.004$); the continuous per-mile result agrees in sign ($p{<}0.001$).
\item A formal pre-trends test finds mean reversion in the raw switcher comparison; a baseline-rate control fixes it.
\item The continuous result, unlike the switcher result, survives state and county fixed effects, though both lose significance under the sparser state-by-transition interaction, read as low-power rather than a reversal.
\item A two-way fixed-effects vs. Callaway--Sant'Anna estimator comparison shows the stacked design sits between a biased-upward TWFE estimate and an underpowered CS estimate whose own pre-trends corroborate our mean-reversion diagnosis; adding our baseline-rate covariate to CS directly does not fix this and instead fails specifically on the one cohort our own diagnostics already flag as imbalanced, at all three distance thresholds checked.
\item A dispensary-level robustness check, using real geocoded license-opening dates in place of state-legal-county distance, reproduces the same sign and a similar magnitude on a smaller, 9-state-coverage sample.
\item Using already-near counties as an alternative control group on both treatment-timing definitions, the switcher coefficient is never significant against the state-legal-county measure and only survives without controls against the dispensary measure -- a real limitation, reported transparently.
\item A medical-legalization placebo check is null; the effect does not compound in years after switching.
\item Implies $\sim$755 averted deaths, $\sim$22,500 life-years saved -- illustrative, not confirmed.
\end{itemize}

\vspace{0.4cm}

% ============================================================
% 1. INTRODUCTION
% ============================================================
\section{Introduction}

Does proximity to a legal recreational-marijuana market change
drug-overdose mortality for residents who never cross the border
themselves? The evidence on marijuana-opioid substitution is
unsettled: \citet{Bachhuber2014} found state medical-marijuana laws
associated with fewer opioid deaths, \citet{Shover2019} showed the
association reverses sign with more years of data, and
\citet{PowellPaculaJacobson2018} found any effect specific to
dispensary-based designs. \citet{MathurRuhm2022} conclude retail
dispensaries are, if anything, associated with \emph{higher} opioid
mortality. Most recently, \citet{Denkyirah2026} use a state-level
two-way-fixed-effects design for \emph{recreational} legalization and
find a decline of roughly 3.5 deaths per 100,000 -- a magnitude we
return to below.

Every one of these papers identifies off policy \emph{timing}. None
uses \emph{geography}: whether proximity to an external legal market
predicts a non-adopting county's own overdose trajectory, in the
spirit of \citet{HaoCowan2020}'s cross-border design for
marijuana-possession arrests. We apply that design to overdose
mortality using both a binary switcher indicator and a continuous,
per-mile distance measure, built as a stacked multi-period panel that
uses all 24 states' actual legalization dates rather than collapsing a
multi-year staggered rollout into two snapshots. An earlier,
two-snapshot version of this design found a null, sign-flipped
continuous estimate that we could not resolve with a bandwidth-based
fix; the multi-period panel resolves it directly. The multi-period
panel also makes possible three checks unavailable in the two-snapshot
version: a formal pre-trends test, a persistence (post-switch
dynamics) check, and a county fixed-effects robustness test -- all
reported below, together with a placebo-style check that swaps in
medical- for recreational-marijuana legalization timing, and a further
robustness check that replaces our county-centroid distance measure
with distance to real, geocoded, individually-dated dispensaries where
such records are publicly available.

% ============================================================
% 2. DATA AND EMPIRICAL STRATEGY
% ============================================================
\section{Data and empirical strategy}

We use 2,331 counties in the 31 U.S. states with no
recreational-marijuana retail market as of our latest
snapshot.\footnote{The companion project's original 32-state
never-legalizing panel included Maine; Maine's retail market opened
October 2020, before even our earliest snapshot, so its counties are
excluded here.} Each county's distance to the nearest legal-market
county is computed via Haversine distance between county centroids
(Census Gazetteer), with ``legal-market'' defined by first retail
sales date, for every December 2020--2025 -- the full window CDC
publishes county-level 12-month-ending overdose data for. Rather than
compare two fixed snapshots, we use all 24 states' actual
retail-opening dates and build a stacked panel of year-over-year
transitions. In transition $t-1\to t$, a county is a ``new switcher''
if its distance to the nearest legal-market county crossed from at
least 50 miles to below 50; a ``control'' if it stayed at least 50
miles away at both endpoints. Once a county switches it is dropped
from every later transition, avoiding the already-treated-unit-as-
control bias documented in the staggered difference-in-differences
literature \citep{GoodmanBacon2021}. Counties within 50 miles as of
the December-2020 baseline (116 of 2,331) are ``always-near'' and
dropped throughout.

\emph{Why 50 miles.} This radius is calibrated qualitatively rather
than to a marijuana-specific travel-distance estimate: the
cross-border consumption literature establishes that cigarette- and
alcohol-tax avoidance is concentrated near a price or legal-status
border and decays with distance to it \citep{Lovenheim2008,
BeattyLarsenSommervoll2009}, and \citet{HaoCowan2020}'s own
cross-border marijuana-spillover design uses literal county-border
adjacency -- a tighter concept than a 50-mile radius -- so our
threshold can be read as a deliberately wider generalization of that
same border-proximity logic to a continuous distance measure, not an
unrelated pick; we stop short of claiming either paper pins down 50
miles specifically, since neither reports a directly comparable
calibrated radius. The choice is also bounded by a real data
constraint: VSRR's small-cell suppression (discussed below)
disproportionately removes sparsely-populated counties, so narrower
radii shrink an already-suppressed sample further, compounding rather
than merely reducing statistical power. Table~\ref{tab:threshcomposition}
shows this tradeoff directly: the number of exposed counties (and the
population share they represent) grows roughly monotonically from 126
counties (6.7\% of the sample population) at 20 miles to 852 counties
(33.7\% of the population) at 100 miles, with 50 miles -- 403 exposed
counties, 16.9\% of the population -- sitting well clear of the
thinnest, least reliable end of that range without classifying a third
of the sample as ``exposed.'' We do not claim 50 miles corresponds to
a discovered discontinuity in overdose mortality itself: a direct
search for a break in mortality as a function of distance finds no
visually distinct kink at 50, 70, or 100 miles specifically
(Figures~\ref{fig:distscan}--\ref{fig:rdplot50}, Section~\ref{sec:appendix}).
The threshold should therefore be read as a designed cutoff,
cross-validated by its downstream stability -- estimates are
qualitatively similar when the identical specification is re-estimated
at 70 and 100 miles (Table~\ref{tab:thresh-main}) -- rather than as an
empirically discovered break point.

\begin{table}[htbp]
\centering
\caption{Sample composition at alternative exposure-distance thresholds}
\label{tab:threshcomposition}
\resizebox{\textwidth}{!}{%
\begin{tabular}{@{}lrrrrr@{}}
\toprule
Threshold & Exposed & Control & \% counties & Population & \% population \\
 & counties & counties & exposed & exposed & exposed \\
\midrule
20 miles & 126 & 2{,}200 & 5.4\% & 12.0M & 6.7\% \\
25 miles & 187 & 2{,}139 & 8.0\% & 16.4M & 9.1\% \\
40 miles & 317 & 2{,}009 & 13.6\% & 26.7M & 14.8\% \\
\textbf{50 miles (primary)} & \textbf{403} & \textbf{1{,}923} & \textbf{17.3\%} & \textbf{30.4M} & \textbf{16.9\%} \\
\textbf{70 miles (robustness)} & \textbf{585} & \textbf{1{,}741} & \textbf{25.2\%} & \textbf{42.6M} & \textbf{23.7\%} \\
\textbf{100 miles (robustness)} & \textbf{852} & \textbf{1{,}474} & \textbf{36.6\%} & \textbf{60.7M} & \textbf{33.7\%} \\
\bottomrule
\end{tabular}%
}
\vspace{2pt}
\begin{minipage}{0.92\textwidth}
\footnotesize\emph{Notes:} December 2023 snapshot, 2,326 counties in
the 31-state never-legal universe (Hawaii excluded; its distance to
the nearest legal-market county is not comparable). ``Exposed''
counties are within the stated distance of the nearest legal-market
county; total sample population is 180.0 million. Bold rows are the
three thresholds used elsewhere in this paper (50 miles primary, 70
and 100 miles as robustness checks, Table~\ref{tab:thresh-main}); 20
and 25 miles are shown for comparison only and are not used in any
regression -- at these radii the switcher pool becomes too small, once
combined with VSRR's small-cell suppression, to support reliable
estimation.
\end{minipage}
\end{table}

The unit of observation is a county-transition pair $(i,t)$: county
$i$ contributes one row for each one-year transition
$t\in\{\text{2020--21},\ldots,\text{2024--25}\}$ it is ``at risk'' of
switching in -- every transition up to and including the one in which
it switches, if it ever does, after which it is dropped, as just
described. $\Delta Y_{it}$ is the change in county $i$'s overdose
death rate per 100,000 between the start and end of transition $t$.
$\text{Switcher}_{it}$ is a binary indicator equal to 1 if and only if
county $i$'s distance to the nearest legal-market county crossed from
at least 50 miles to below 50 miles \emph{during} transition $t$, so a
given county contributes at most one $\text{Switcher}_{it}=1$
observation across the whole panel; it equals 0 for every transition a
county has not yet switched in, including every transition for
counties that never switch. $\Delta\text{Dist}_{it}$ is the raw change
in miles to the nearest legal-market county over transition $t$,
defined for every county-transition row regardless of switcher status
(control counties can still move closer or farther without crossing
the threshold). $\delta_t$ denotes transition (period) fixed effects,
common to all counties observed in a given transition. Standard errors
are clustered by state, the level at which most of the underlying
legal-market-timing variation occurs. We pool all (county, transition)
observations into one stacked regression
\citep{CengizDubeLindnerZipperer2019}:
\begin{equation}
  \Delta Y_{it} = \alpha + \gamma \cdot \text{Switcher}_{it} + \delta_t + \varepsilon_{it}, \qquad
  \Delta Y_{it} = \alpha + \beta \cdot \Delta\text{Dist}_{it} + \delta_t + \varepsilon_{it}.
  \label{eq:main}
\end{equation}
Equation~\ref{eq:main} imposes one switcher effect $\gamma$ and one
distance effect $\beta$ common to all five transitions. We also
estimate a more general version that lets both vary freely by period,
via a full interaction with the transition dummies:
\begin{equation}
  \Delta Y_{it} = \alpha + \sum_t \gamma_t \cdot (\text{Switcher}_{it} \times \mathbb{1}[t])
                        + \sum_t \beta_t \cdot (\Delta\text{Dist}_{it} \times \mathbb{1}[t])
                        + \delta_t + \varepsilon_{it},
  \label{eq:period}
\end{equation}
which nests Equation~\ref{eq:main} as the special case
$\gamma_t=\gamma$ and $\beta_t=\beta$ for all $t$;
Section~\ref{sec:periodeffects} reports the period-specific
coefficients and tests whether pooling is appropriate.

As of December 2020, ten states had an open retail market; by December
2025, 23 did (Figure~\ref{fig:rollout}). Virginia has never opened a
general retail market and so never enters the legal-market set.
The second source is county-level drug-overdose mortality. This comes
from the National Center for Health Statistics (NCHS) -- the CDC's
vital-statistics agency -- through its Vital Statistics Rapid Release
(VSRR) program, which publishes provisional county-level death counts
drawn from the same current flow of death-certificate records that
feeds NCHS's final annual mortality files, but on a much faster cycle
and before all cause-of-death investigations are complete
\citep{AhmadVSRR, RossenNowcasting2021}; CDC WONDER is not reachable
from our environment. We use these raw provisional counts rather than
NCHS's own smoothed, model-based small-area alternative, so that a
genuine, sudden, localized change at a legal-market border is not
averaged away by construction -- the same reasoning that leads other
studies in this area to prefer raw counts, with their own thresholds
for handling small-number instability, over smoothed county estimates
\citep{SabaEtAl2025}.

Two features of this series bind directly on our design. First,
NCHS's confidentiality standard -- which applies to VSRR just as it
does to the final files -- suppresses any county-transition cell with
1--9 deaths at either endpoint, a real limitation of working with
county-level rather than state-level mortality counts. Second,
county-level detail in VSRR is available only from 2020 onward
\citep{AhmadVSRR}, which is why this paper's panel begins in 2020
rather than earlier, and cannot speak to the earlier 2014--2018
legalization wave (Colorado, Washington, Oregon, Alaska, California).
Suppression leaves 4,259 of 10,313 potential observations (177 switch
events).

As a placebo-style check, we rebuild the identical stacked design
using each state's medical-marijuana access date (first regulated
medical dispensary sales, compiled for 39 states) in place of the
recreational-retail date, restricting the county universe to the 17 of
31 states not yet medically legal as of the December-2020 baseline.

\begin{figure}[htbp]
  \centering
  \includegraphics[width=\textwidth]{figures/fig1_rollout_maps_multiperiod.png}
  \caption{Legal-market rollout (by first retail sales date) and 50-mile treatment frontier, 2020--2025.}
  \label{fig:rollout}
\end{figure}

% ============================================================
% 3. RESULTS
% ============================================================
\section{Results}

\begin{table}[htbp]
\centering
\caption{Effect on drug-overdose mortality, preferred specification}
\label{tab:main}
\begin{tabular}{@{}lcc@{}}
\toprule
 & Switcher (binary) & Distance, per mile \\
\midrule
Overdose rate per 100k & $-2.320^{***}$ & $+0.0122^{***}$ \\
SE (state-clustered) & (0.800) & (0.0034) \\
SE (county-clustered) & (0.776) & (0.0031) \\
$p$ (state-clustered) & 0.004 & $<0.001$ \\
$p$ (county-clustered) & 0.003 & $<0.001$ \\
\bottomrule
\end{tabular}
\vspace{2pt}
\begin{minipage}{0.85\textwidth}
\footnotesize\emph{Notes:} Outcome is the transition-level change in
overdose deaths per 100{,}000. Both columns include transition fixed
effects and a control for each county's baseline (2020) overdose-rate
level (Section~\ref{sec:pretrends}). $N=3{,}794$, 31 state clusters,
1{,}249 county clusters. The distance coefficient is signed as the
effect of an \emph{increase} in distance, so it agrees in direction
with the (negative) switcher coefficient. Stars reflect the
state-clustered $p$-value, our primary (and more conservative)
inference throughout the paper. $^{***}p<0.01$.
\end{minipage}
\end{table}

\begin{figure}[htbp]
  \centering
  \includegraphics[width=0.85\textwidth]{figures/fig2_overdose_coefplot_multiperiod.png}
  \caption{Effect on overdose mortality, recreational (primary) vs. medical (placebo) design (95\% CI).}
  \label{fig:coef}
\end{figure}

Switcher counties see overdose deaths fall 2.32 more per 100,000 than
controls in the same transition ($p=0.004$; Table~\ref{tab:main}) --
somewhat below \citet{Denkyirah2026}'s state-level estimate of roughly
3.5, though the same direction and order of magnitude. The continuous
per-mile specification agrees in direction and is itself significant
($p<0.001$).

\subsection{Pre-trends check and robustness to fixed effects}
\label{sec:pretrends}

Because eventual-switcher counties contribute observations to every
transition before they switch, we can test whether their overdose
trend was already diverging from never-switching counties
\emph{before} switching. The raw version of this test fails: leads on
event-time indicators (one and two-plus transitions before switching)
are large and jointly significant ($p<0.001$), comparable in magnitude
to the switch-transition effect itself, even though transition fixed
effects are already present. We traced this to mean reversion --
counties with unusually high 2020 overdose rates mechanically decline
more -- and adding each county's baseline (2020) rate as a control
resolves it (joint F-test on both leads, $p=0.318$; Figure~\ref{fig:pretrends}).
We also checked Census-division-by-transition fixed effects (a
geographic-clustering explanation); this also passes the joint test on
this panel but leaves the switch coefficient less precisely estimated
and proved sensitive to exact sample construction across our
diagnostic checks, so we prefer the baseline-rate control
(Table~\ref{tab:main}).

\begin{figure}[htbp]
  \centering
  \includegraphics[width=0.85\textwidth]{figures/fig5_eventstudy_pretrends.png}
  \caption{Pre-trends check: raw vs. mean-reversion-adjusted event-study coefficients.}
  \label{fig:pretrends}
\end{figure}

\begin{table}[htbp]
\centering
\caption{Robustness of main coefficients to fixed-effects controls}
\label{tab:fe}
\resizebox{\textwidth}{!}{%
\begin{tabular}{@{}lccccc@{}}
\toprule
 & \multicolumn{3}{c}{Switcher (binary)} & \multicolumn{2}{c}{Distance (per mile)} \\
Specification & Coef. & $p$ (state) & $p$ (county) & Coef. & \multicolumn{1}{c}{$p$ (state) / $p$ (county)} \\
\midrule
No controls & $-3.151$ & 0.005 & $<0.001$ & $+0.0128$ & $<0.001$ / $<0.001$ \\
+ Baseline overdose-rate level (preferred) & $-2.320^{***}$ & 0.004 & 0.003 & $+0.0122$ & $<0.001$ / $<0.001$ \\
+ State fixed effects & $-1.646$ & 0.109 & 0.062 & $+0.0125$ & $<0.001$ / $<0.001$ \\
+ Census division $\times$ transition FE & $-1.519$ & 0.109 & 0.084 & $+0.0084^{***}$ & 0.004 / 0.011 \\
+ State $\times$ transition FE & $+2.015$ & 0.168 & 0.141 & $+0.0026$ & 0.606 / 0.572 \\
+ County fixed effects & $-0.999$ & 0.532 & 0.463 & $+0.0122^{**}$ & 0.023 / 0.009 \\
\bottomrule
\end{tabular}%
}
\vspace{2pt}
\begin{minipage}{0.95\textwidth}
\footnotesize\emph{Notes:} Each row is a separate regression on the
same $N=3{,}794$ sample; 31 state clusters, 1{,}249 county clusters.
The census-division and state $\times$ transition rows also include
the baseline-rate control. County fixed effects use only
within-county variation across a given county's (at most five)
transition-rows -- 225 of 1,249 counties in this sample have only one
transition-row and so contribute no within-county variation; clustering
by county on this row groups observations at the same level as the
fixed effect itself. Stars reflect the state-clustered $p$-value.
$^{**}p<0.05$, $^{***}p<0.01$. \emph{Correction:} the ``+ State fixed
effects'' row previously printed $-1.411$ ($p{=}0.144$) and $+0.0129$
($p{<}0.001$); those numbers did not match a fresh re-estimation on
the current panel (every other row did) and have been corrected here
to $-1.646$ / $+0.0125$. The qualitative conclusion --  the switcher
coefficient is insignificant under state FE while the distance
coefficient remains significant -- is unchanged.
\end{minipage}
\end{table}

Table~\ref{tab:fe} shows the two specifications respond very
differently to fixed-effects controls. The continuous per-mile
coefficient is essentially unmoved by state fixed effects, census
division-by-transition fixed effects, county fixed effects, \emph{or}
the baseline-rate control, and remains significant throughout --
except under state-by-transition fixed effects, where it falls to
$+0.0026$ ($p=0.606$) and the switcher coefficient flips sign to
$+2.015$ ($p=0.168$, also insignificant). We read this as the sparsest,
most overfit specification in the table rather than a genuine
reversal: it spreads 3,794 observations across up to 155
state-by-transition cells (31 states $\times$ 5 transitions), many
containing only one or two counties, the same low-power pattern
already documented for the analogous rows in
Tables~\ref{tab:dispensary} and~\ref{tab:alwaysnear}. Census
division-by-transition fixed effects fragment the sample far less (9
divisions $\times$ 5 transitions, at most 45 cells) and the continuous
coefficient stays significant there ($p=0.004$), though the switcher
coefficient does not ($-1.519$, $p=0.109$), matching the qualitative
pattern from the rest of this table. County fixed
effects -- which use only within-county variation but do not fragment
the sample into nearly as many near-empty cells -- remain, in our
judgment, the more informative of the two demanding tests, and the
continuous specification survives that one. The binary switcher
coefficient survives the baseline-rate control but loses significance
under state fixed effects and further under county fixed effects,
indicating a meaningful share of its identifying variation comes from
\emph{which} counties or states happen to border a soon-to-open
market, not from finer within-unit proximity differences. We therefore
treat the continuous per-mile result as this paper's primary
quantitative claim, with the switcher result as corroborating but
comparatively more fragile.

\subsection{Estimator robustness: two-way fixed effects and Callaway--Sant'Anna}
\label{sec:twfecs}

The stacked-cohort estimator in Equation~\ref{eq:main} is one
particular way to handle staggered treatment timing; it is not the
only one, and the staggered difference-in-differences literature has
shown that the more familiar two-way fixed-effects (TWFE) regression
can be badly biased under staggered adoption with heterogeneous
treatment effects \citep{GoodmanBacon2021}. We check whether our
result is an artifact of the stacking approach itself by re-estimating
it two more ways on the same underlying treatment definition (state
legal-market distance, county crosses below 50 miles): a canonical
TWFE regression, and the doubly-robust group-time estimator of
\citet{CallawaySantAnna2021}.

Both require a genuine county-by-year panel with an \emph{absorbing}
treatment indicator (once a county crosses the 50-mile threshold it
stays ``treated'' in every subsequent year, rather than being dropped,
as the stacked design does) and, unlike the stacked design's
transition-by-transition suppression handling, a county with any
suppressed year cannot contribute at all. This is a materially
stricter data requirement: of the same 2,215 never-legal, not-always-near
counties used elsewhere, only 614 have non-suppressed overdose data in
every one of the six panel years (2020--2025), giving a balanced panel
of 3,684 county-year observations with three identified treatment
cohorts (39 counties first crossing in 2022, 41 in 2023, 65 in 2024;
counts differ slightly from the switch counts elsewhere because this
panel requires complete data across \emph{all six} years, not just the
two years spanning a given transition).

The TWFE regression is $Y_{it} = \alpha_i + \lambda_t + \beta \cdot
\text{Treated}_{it} + \varepsilon_{it}$, county and year fixed effects,
clustered by state. The \citet{CallawaySantAnna2021} estimator computes
group-time average treatment effects $ATT(g,t)$ for each treatment
cohort $g$ against the never-treated counties, then aggregates first to
an overall average effect and separately to an event-study (dynamic)
profile by time since treatment, using the multiplier bootstrap for
inference.

\begin{table}[htbp]
\centering
\caption{Estimator comparison: TWFE and Callaway--Sant'Anna vs. the stacked-cohort design}
\label{tab:twfecs}
\resizebox{\textwidth}{!}{%
\begin{tabular}{@{}lcccccc@{}}
\toprule
 & & \multicolumn{2}{c}{State-clustered} & \multicolumn{2}{c}{County-clustered} & \\
Estimator & Coef. & SE & $p$ & SE & $p$ & Clusters (st./cty.) \\
\midrule
Stacked-cohort switcher (Table~\ref{tab:main}, preferred spec) & $-2.320$ & $0.800$ & $0.004$ & $0.776$ & $0.003$ & 31 / 1{,}249 \\
TWFE (county + year FE) & $-5.727$ & $2.338$ & $0.014$ & $1.197$ & $<0.001$ & 31 / 614 \\
Callaway--Sant'Anna, overall ATT, never-treated controls & $-1.905$ & $1.929$ & $0.323$ & $1.038$ & $0.067$ & 31 / 614 \\
\quad + baseline-rate covariate ($x$formla) & $+5.992$ & $3.163$ & $0.058$ & $4.116$ & $0.145$ & 31 / 614 \\
\bottomrule
\end{tabular}%
}
\vspace{2pt}
\begin{minipage}{0.95\textwidth}
\footnotesize\emph{Notes:} TWFE and Callaway--Sant'Anna use the
balanced 614-county panel described above ($n=3{,}684$ county-years);
the stacked-cohort row is reproduced from Table~\ref{tab:main} on its
own, larger sample for reference and is not directly comparable
observation-for-observation. All three use the same underlying
treatment definition (distance to nearest legal-market county crossing
50 miles). County-clustered SE for TWFE clusters at the same level as
its own county fixed effect. The covariate-adjusted Callaway--Sant'Anna
row conditions the doubly-robust estimator on each county's baseline
(2020) overdose-rate level, the same variable used as the stacked
design's mean-reversion control (Section~\ref{sec:pretrends}); see
discussion below for why this row should not be read as a stable
estimate.
\end{minipage}
\end{table}

Table~\ref{tab:twfecs} shows the three estimators do not agree closely
with each other, and we think that disagreement is itself informative.
TWFE returns a coefficient more than double the stacked design's
preferred estimate ($-5.73$ vs.\ $-2.32$) -- consistent with the
well-documented direction of TWFE bias under staggered adoption, where
already-treated units act as ``forbidden comparisons'' for later-treated
units and can inflate the estimated effect. Callaway--Sant'Anna, by
contrast, returns a smaller, statistically insignificant overall ATT
($-1.91$, $p=0.32$) on the smaller balanced sample. Examining
Callaway--Sant'Anna's own event-study decomposition
(Figure~\ref{fig:cs_eventstudy}) shows why: the pre-treatment
coefficients at $t=-1$ and $t=-2$ are themselves significant
($p=0.031$ and $p<0.001$), the same mean-reversion-driven pre-trend
contamination diagnosed in Section~\ref{sec:pretrends} -- except this
estimator, run in its default form, has no analog of our baseline-rate
covariate to correct for it. Rather than treating either off-the-shelf
estimator as a referee on the stacked design, we read this as
corroborating the pre-trends finding from an entirely independent
estimator and confirming why we build a stacked cohort design with an
explicit mean-reversion control rather than applying either canonical
staggered-adoption estimator directly.

A natural follow-up is to give Callaway--Sant'Anna the same covariate
we use, via its built-in covariate-adjustment option (a doubly-robust
correction conditioning the group-time ATTs on the county's baseline
overdose-rate level). If the pre-trend problem is exactly what the raw
CS event-study suggests, this should shrink the pre-treatment
coefficients toward zero and move the overall ATT toward the stacked
design's $-2.32$. It does not: the covariate-adjusted overall ATT is
$+5.99$ ($se=3.16$, $p=0.058$) -- the wrong sign, and the pre-treatment
coefficients remain significant, now at $t=-3$ and $t=-1$
($p=0.0007$ and $p=0.040$) rather than $t=-2$ and $t=-1$, with the
whole event-study profile swinging between $-11.4$ and $+15.1$ across
adjacent periods (Figure~\ref{fig:cs_eventstudy}, right panel).

Tracing this to its source, the instability is not general -- it is
entirely a property of one cohort. The group-time cells for counties
first treated in 2024 (the Ohio-border transition already flagged in
Section~\ref{sec:periodeffects} for its severe baseline-rate imbalance
against its control pool, 53.4 vs.\ 26.1) swing wildly even in
placebo/pre-treatment periods once the covariate is added
($-11.4$, $+11.6$, $-13.0$ across three pre-treatment years); the 2022
and 2023 cohorts' cells stay small and stable throughout. Re-estimating
with the 2024 cohort excluded gives a small, stable, insignificant
overall ATT ($+0.77$, $se=0.76$, $p=0.31$). At the 100-mile threshold
the same cohort's imbalance is severe enough that the doubly-robust
estimator cannot compute its cells at all (every group-time ATT for
that cohort is missing); dropping it automatically leaves a small,
correctly-signed, insignificant estimate ($-0.41$, $p=0.66$). This
pattern -- sign-flipping or non-estimable exactly where covariate
overlap between a treated cohort and its controls is worst -- is the
well-documented failure mode of propensity/outcome-regression-based
covariate adjustment under poor overlap, and it holds at all three
distance thresholds we check (Table~\ref{tab:thresh-twfecs}). We read
this as a second, independent confirmation of the same cohort-specific
imbalance Section~\ref{sec:periodeffects} already diagnoses, this time
surfacing as an outright estimation failure rather than a coefficient
requiring investigation -- not as evidence against covariate adjustment
in general, but as a concrete illustration of why the stacked design's
simple additive baseline-rate control, paired with the explicit
cohort-level diagnostics in Section~\ref{sec:periodeffects}, is more
robust here than handing the correction to an off-the-shelf
covariate-balancing routine.

\begin{figure}[htbp]
  \centering
  \includegraphics[width=\textwidth]{figures/figH_cs_eventstudy.png}
  \caption{Callaway--Sant'Anna event-study coefficients by time since treatment (never-treated control group). Left: no covariates. Right: + baseline-rate covariate.}
  \label{fig:cs_eventstudy}
\end{figure}

\subsection{Robustness: real dispensary-opening dates in place of state legalization dates}
\label{sec:dispensary}

Every result so far defines ``legal-market distance'' at the state
level: a county's distance to the nearest county in a state whose
retail market is open, as of that state's own first-retail-sales date.
This is a coarse proxy for the thing that should actually matter to a
nearby resident -- physical distance to a place that sells recreational
marijuana. If the state-level measure is systematically off (e.g.\ a
county classified as ``near'' the day its neighboring state legalizes,
even though no dispensary within 50 miles has actually opened yet), it
could bias the estimated speed, though not necessarily the existence,
of the effect. We test this directly by rebuilding the identical
stacked switcher/continuous design with distance measured to the
nearest \emph{individual, currently operating} dispensary rather than
to the nearest legal-market county.

\emph{Data.} We collected official dispensary-license rosters directly
from state cannabis-regulatory agencies (not third-party aggregators)
for every state where a usable, dated, machine-readable roster is
publicly available: Ohio, New York, Massachusetts, Oregon, Illinois,
Maine, and California carry a genuine per-dispensary license-issuance
or first-retail-sale date; Washington and Colorado's own license files
carry only a renewal/expiration date, not an opening date, but both
states legalized recreational retail in 2014 -- years before our
2020--2025 panel begins -- so their dispensaries are coded as already
open throughout the panel rather than assigned a misleading proxy
date. This covers 9 of the 24 legal states; the remaining 15 either
have no public per-license date field (confirmed for Arizona,
Connecticut, Missouri, Montana, New Jersey, Delaware, Nevada, and
Vermont) or were not pursued further given time constraints (Arizona's
own market is small; several others are mid-sized). Records were
geocoded to latitude/longitude via the Census Bureau's public batch
geocoder (98.3\% match rate combining an exact street-address match for
91.7\% of records and a ZIP-centroid fallback for the remainder;
California's own regulator API already returns geocoded coordinates).
After dropping records with no usable coordinate or opening date, 5,178
of 5,825 collected dispensary licenses enter the distance calculation.

This is a genuine limitation, not a footnote to skip past: a
never-legal county whose true nearest dispensary sits in one of the 15
uncovered states is measured here as farther from a dispensary than it
actually is, which pushes that county toward the ``control'' group even
if it should be a treated ``switcher.'' This attenuates the estimated
effect and shrinks the usable switcher count rather than biasing the
sign, but it means this check should be read as confirming
\emph{direction}, not as a like-for-like replication of the main
design's precision.

\emph{Design.} Everything else is identical to Section~2: the same
2,331-county never-legal universe, the same 50-mile switcher/control
threshold, the same drop-after-switching stacked-cohort construction,
and the same CDC VSRR overdose outcome. The only change is that a
county's distance in December of year $y$ is now computed to the
nearest dispensary with an issuance date on or before December 31 of
year $y$ (Washington and Colorado dispensaries are treated as
open in every year). Fifty-three counties are ``always-near'' a real
dispensary as of the 2020 baseline under this measure -- notably fewer
than the 116 dropped under the state-level measure, since a specific
dispensary address is on average farther away than the nearest county
in \emph{any} legal state. The resulting stacked panel has 240 raw
switch events; after the same NCHS small-cell suppression rule, 84
switch events and $N=3{,}941$ county-transition observations remain.

Table~\ref{tab:dispensary} reports the continuous distance coefficient
-- our primary measure throughout this paper -- across the same
progression of fixed-effects controls as Table~\ref{tab:fe}, plus two
additional interacted specifications (state-by-transition and
census-division-by-transition fixed effects) that this smaller sample
can still support.\footnote{We could not extend this fixed-effects
menu to census-tract or calendar-season fixed effects, both requested
as robustness dimensions. Neither is estimable with the data actually
available: CDC's overdose-mortality series (VSRR and, at finer detail,
NCHS restricted files) is published only at the county level -- no
public source reports drug-overdose deaths by census tract, for
confidentiality reasons consistent with NCHS's own county-level
small-cell suppression rule already binding in this design. VSRR does
report monthly, but each monthly value is itself a 12-month
\emph{rolling} total (a moving average constructed specifically to
smooth month-to-month noise in provisional data), so it contains no
genuine within-year seasonal signal to identify a season effect
against -- a season dummy fit to a 12-month rolling series is fitting
noise, not seasonality. We use the richest fixed-effects menu the data
can actually support: county, state, census division, and
time (transition-period) fixed effects, and interactions among them.}

\begin{table}[htbp]
\centering
\caption{Dispensary-level robustness check: distance to nearest operating dispensary}
\label{tab:dispensary}
\resizebox{\textwidth}{!}{%
\begin{tabular}{@{}lcccccc@{}}
\toprule
 & \multicolumn{3}{c}{Distance (per mile)} & \multicolumn{3}{c}{Switcher (binary)} \\
Specification & Coef. & $p$ (state) & $p$ (county) & Coef. & $p$ (state) & $p$ (county) \\
\midrule
No controls (transition FE only) & $+0.0230$ & 0.068 & $<0.001$ & $-5.040^{***}$ & $<0.001$ & $<0.001$ \\
+ Baseline (2020) overdose rate & $+0.0173$ & 0.082 & $<0.001$ & $-2.715^{**}$ & 0.032 & 0.046 \\
+ State FE & $+0.0183$ & 0.099 & $<0.001$ & $-1.936$ & 0.178 & 0.199 \\
+ Census division $\times$ transition FE & $+0.0171$ & 0.108 & $<0.001$ & $-2.034^{**}$ & 0.019 & 0.222 \\
+ State $\times$ transition FE & $+0.0051$ & 0.360 & 0.185 & $+3.632$ & 0.183 & 0.184 \\
+ County FE & $+0.0216$ & 0.124 & $<0.001$ & $-2.041$ & 0.275 & 0.309 \\
\bottomrule
\end{tabular}%
}
\vspace{2pt}
\begin{minipage}{0.95\textwidth}
\footnotesize\emph{Notes:} $N=3{,}941$ throughout; 24 state clusters,
1{,}261 county clusters. All specifications include the baseline-rate
control except the first row. Distance is signed identically to
Table~\ref{tab:fe} (effect of an \emph{increase} in distance to the
nearest operating dispensary). Stars reflect the state-clustered
$p$-value. The continuous coefficient is far more significant under
county clustering throughout (every row $p<0.001$ except the
state-by-transition row), the same pattern as Table~\ref{tab:fe}.
$^{**}p<0.05$, $^{***}p<0.01$.
\end{minipage}
\end{table}

\begin{figure}[htbp]
  \centering
  \includegraphics[width=0.85\textwidth]{figures/figG_dispensary_vs_main_coefplot.png}
  \caption{Continuous distance coefficient: dispensary-opening-date design vs. main (state-legalization-date) design, by fixed-effects specification.}
  \label{fig:dispensary}
\end{figure}

The continuous coefficient stays positive -- the same sign as the main
design throughout -- and of a broadly similar magnitude (+0.005 to
+0.023 per mile, against the main design's +0.012 to +0.013;
Figure~\ref{fig:dispensary}), and the switcher coefficient reproduces
the main design's own qualitative finding that the binary measure is
the more fragile of the two: it survives the baseline-rate control and
division-by-transition fixed effects, but loses significance under
state fixed effects, county fixed effects, and (with the ``wrong''
sign, though far from significant) state-by-transition fixed effects.
None of the continuous coefficients here reach conventional
significance at better than the 10\% level, and state-by-transition
fixed effects null it out entirely. Given how thin this specification
spreads the data -- 31 states $\times$ 5 transitions, many cells
containing only one or two counties -- we read the state-by-transition
result as a low-power, likely overfit specification rather than a
genuine contradiction, but we report it rather than omit it.

On balance, this check supports the main design's qualitative
conclusion (proximity to an actual, operating retail market is
associated with lower overdose mortality) using an entirely
independent, arguably more accurate treatment-timing measure, while
underscoring that the main design's stronger statistical significance
partly reflects its full 24-state national coverage; a 9-state,
partial-coverage version of the same underlying idea is directionally
consistent but not independently conclusive on its own.

\subsection{Robustness: always-near counties as an alternative control group}
\label{sec:alwaysnear}

Every specification so far compares switcher counties against
far/never-treated counties -- counties that stayed at least 50 miles
from any legal-market county throughout. This check instead uses
counties that were \emph{already} near a legal market as an
alternative control group, and does so on \emph{both} of this paper's
treatment-timing definitions: Panel A uses the state-legal-county
distance measure of Section~2 (switcher/always-near defined by
distance to the nearest legal-market county); Panel B uses the
dispensary-level measure of Section~\ref{sec:dispensary} (switcher/
always-near defined by distance to an \emph{actually operating}
dispensary). In both panels, ``always-near'' counties are the ones that
were already within 50 miles of the relevant target as of the
December-2020 baseline -- 116 counties for Panel A (near Illinois,
Massachusetts, Michigan, Nevada, and the first-wave 2014--2018 states),
53 for Panel B (concentrated in Iowa, Idaho, Indiana, Kentucky,
Nebraska, New Hampshire, Utah, Wisconsin, and Wyoming -- e.g., Cheyenne,
WY, sits within 50 miles of dispensaries operating in the Fort Collins,
CO area). Switcher status is otherwise unchanged in each panel: every
county that ever switches contributes its full history of at-risk
transitions (both its pre-switch $\text{Switcher}=0$ years and its
switch-transition $\text{Switcher}=1$ year), while always-near counties
enter every one of the five transitions as $\text{Switcher}_{it}=0$,
since their distance to the relevant target never crosses the 50-mile
line during our panel -- it was already under it before the panel
began. Ever-switcher counties' own pre-switch years must be included
this way, not dropped down to only the switch-moment row, for a county
fixed effect to have any within-county variation to work with at
all.\footnote{We could not add a census-tract fixed effect here, for
the same reason given in Section~\ref{sec:dispensary}'s footnote: CDC's
overdose-mortality data, including the restricted NCHS files, is
published only at the county level, so no tract identifier exists
anywhere in the outcome data to fit a tract effect against.}

\begin{table}[htbp]
\centering
\caption{Switcher effect against an always-near control group, both treatment-timing definitions}
\label{tab:alwaysnear}
\begin{tabular}{@{}lcccc@{}}
\toprule
Specification & Coef. & $p$ (state) & $p$ (county) & $N$ \\
\midrule
\multicolumn{5}{@{}l}{\emph{Panel A: state-legal-county distance (Section~2)}} \\
No controls (transition FE only) & $-0.797$ & 0.427 & 0.410 & 858 \\
+ Baseline (2020) overdose rate & $-1.537$ & 0.152 & 0.098 & 798 \\
+ State fixed effects & $-1.489$ & 0.236 & 0.184 & 798 \\
+ Census division $\times$ transition FE & $-1.633$ & 0.339 & 0.173 & 798 \\
+ State $\times$ transition FE & $-1.114$ & 0.650 & 0.489 & 798 \\
+ County fixed effects & $-1.066$ & 0.581 & 0.537 & 798 \\
\addlinespace
\multicolumn{5}{@{}l}{\emph{Panel B: distance to an actually operating dispensary (Section~\ref{sec:dispensary})}} \\
No controls (transition FE only) & $-3.116^{**}$ & 0.037 & 0.177 & 468 \\
+ Baseline (2020) overdose rate & $-0.957$ & 0.457 & 0.694 & 443 \\
+ State fixed effects & $-0.719$ & 0.627 & 0.791 & 443 \\
+ Census division $\times$ transition FE & $-0.130$ & 0.927 & 0.954 & 443 \\
+ State $\times$ transition FE & $+2.495^{**}$ & 0.035 & 0.506 & 443 \\
+ County fixed effects & $-0.554$ & 0.768 & 0.866 & 443 \\
\bottomrule
\end{tabular}
\vspace{2pt}
\begin{minipage}{0.95\textwidth}
\footnotesize\emph{Notes:} Binary switcher specification only in both
panels. Panel A pools all at-risk transition-rows for the 236 counties
that ever switch (177 switch-transition rows plus their own pre-switch
history) with the 116 always-near counties' rows (repeated across all
five transitions), 298 county clusters. Panel B pools the 91 counties
that ever switch (91 switch-transition rows plus their own pre-switch
history) with the 53 always-near counties' rows, 121 county clusters
(468 rows/121 counties $\approx$3.9 rows per county, so county
clustering still pools repeat observations, unlike the single-period
tables below). County fixed effects use only within-county variation;
52 of 298 counties in Panel A and 7 of 121 in Panel B have only one row
and so contribute none. Census-tract fixed effects are not estimable
(see footnote). Stars reflect the state-clustered $p$-value. Note
Panel B's county-clustered SEs are consistently \emph{larger}, not
smaller, than its state-clustered SEs -- the opposite pattern from
every other table in this paper -- reflecting this panel's small,
9-state, unevenly distributed cluster structure rather than a general
finding. $^{**}p<0.05$.
\end{minipage}
\end{table}

In Panel A, none of the six specifications reach conventional
significance, including before any fixed effects are added ($-0.80$,
$p=0.427$); census division-by-transition fixed effects give the
largest point estimate of the panel ($-1.63$) but, like every other
row, are far from significant ($p=0.339$). In Panel B, switcher
counties decline significantly more than always-near counties without
controls ($-3.12$, $p=0.037$), but the two groups start from very
different baselines (switcher counties average a 2020 overdose rate of
52.1 per 100,000 against always-near counties' 27.2), and the
coefficient falls by more than half and loses significance once the
baseline-rate control is applied ($-0.96$, $p=0.457$), weakening
further under state fixed effects ($-0.72$, $p=0.627$) and virtually to
zero under census division-by-transition fixed effects ($-0.13$,
$p=0.927$) and county fixed effects ($-0.55$, $p=0.768$). Panel B's
state-by-transition row is the one exception -- it flips sign and
becomes significant ($+2.50$, $p=0.035$) -- but we read this as a
low-power, likely overfit specification rather than a genuine reversal:
it spreads only 443 observations across up to 45 state-by-transition
cells (9 states $\times$ 5 transitions), the same pattern already
flagged for the analogous row in Table~\ref{tab:dispensary}. We read
both panels honestly rather than explain them away: against a
comparison group that already had market access, most of the raw gap
looks like mean-reversion-driven convergence from a higher starting
point rather than a clean differential trend. It does not overturn the
paper's headline claim -- the continuous specification remains our
primary measure and was not re-tested here (this check is binary-only
in both panels, since a meaningful continuous distance change is not
defined for counties whose distance to the relevant target never
moves) -- but it is a genuine limitation on how far the switcher result
generalizes across different, defensible choices of comparison group.

\subsection{Period-specific effects}
\label{sec:periodeffects}

Equation~\ref{eq:main} assumes the switcher and distance effects are
constant across all five transitions. Estimating Equation~\ref{eq:period}
instead -- letting both vary freely by period -- and restricting to
the three transitions with enough switch events to identify a
period-specific coefficient (2020--2021 and 2024--2025 have zero
switch events between them after the baseline-rate restriction, so
contribute to the period fixed effects but not to their own
switcher/distance coefficient), gives Table~\ref{tab:period}.

\begin{table}[htbp]
\centering
\caption{Period-specific switcher and distance coefficients}
\label{tab:period}
\resizebox{\textwidth}{!}{%
\begin{tabular}{@{}lccccc@{}}
\toprule
 & \multicolumn{2}{c}{Switcher (binary)} & \multicolumn{2}{c}{Distance (per mile)} & \\
Transition & Coef. & $p$ (state) / (county) & Coef. & $p$ (state) / (county) & $n$ switchers \\
\midrule
2021--2022 & $-0.648$ & 0.547 / 0.612 & $+0.0014$ & 0.717 / 0.709 & 23 \\
2022--2023 & $-1.718$ & 0.449 / 0.147 & $-0.0027$ & 0.802 / 0.658 & 62 \\
2023--2024 & $+8.280^{***}$ & 0.008 / 0.022 & $+0.1132^{***}$ & $<0.001$ / $<0.001$ & 75 \\
\bottomrule
\end{tabular}%
}
\vspace{2pt}
\begin{minipage}{0.9\textwidth}
\footnotesize\emph{Notes:} One joint regression, all period
interactions plus period fixed effects and the baseline-rate control;
$N=3{,}794$; 31 state clusters, 1{,}249 county clusters. An F-test
rejects that the three switcher coefficients are equal (state-clustered
$p=0.002$, county-clustered $p=0.038$) and that the three distance
coefficients are equal (state-clustered $p=0.003$, county-clustered
$p<0.001$): the pooled Table~\ref{tab:main} coefficient is not simply
the same effect repeated five times. Stars reflect the state-clustered
$p$-value. $^{***}p<0.01$.
\end{minipage}
\end{table}

The 2023--2024 transition -- which contributes 75 of the panel's 177
switch events, concentrated in Ohio-border states (Ohio, West
Virginia, Pennsylvania, Kentucky, Indiana) -- has a positive,
significant coefficient on both specifications, the opposite sign from
the pooled result. We investigated this rather than reporting it at
face value, and three things resolve it. First, in raw, unadjusted
terms switcher counties in this transition declined \emph{more} than
controls (mean $\Delta$ overdose rate $-14.5$ vs.\ $-8.5$); the
positive sign appears only after the baseline-rate control, because
switchers in this specific transition have roughly double the 2020
baseline rate of controls (53.4 vs.\ 26.1), a far more severe imbalance
than in other transitions. Second, this period's control pool is
geographically mismatched to its switcher pool -- comparing
Appalachian counties against a national control pool concentrated in
the South Atlantic and West South Central divisions -- and adding
Census-division fixed effects shrinks the switcher coefficient to
$+2.79$ ($p=0.253$) and the distance coefficient to $+0.015$
($p=0.335$), both no longer significant. Third, a handful of
extreme-outlier small counties (implausibly large overdose-rate swings
driven by small population denominators) have outsized leverage on
this single-period fit; excluding the top 5\% by Cook's distance
likewise removes significance ($+1.82$, $p=0.166$). Combining division
fixed effects with winsorizing the baseline-rate control at its 95th
percentile gives $+2.33$ ($p=0.304$) for the switcher coefficient and a
smaller, though still marginally significant, $+0.032$ ($p=0.036$) for
distance. We read this as within-period estimation noise from
isolating a geographically concentrated, high-baseline-rate cohort with
few observations, not as evidence against the pooled result: it draws
on switch events spread across many more states and periods than any
single transition, and Table~\ref{tab:fe}'s state and county
fixed-effects checks already show the continuous specification is not
being driven by any one state or region.

\subsection{Medical-legalization placebo check}

Re-running the identical stacked design with medical- in place of
recreational-marijuana timing, both the switcher ($b=+0.51$,
se$=1.04$, $p=0.62$) and continuous ($b=-0.015$, se$=0.008$, $p=0.08$)
estimates are null -- consistent with the recreational-retail effect
reflecting retail access specifically rather than the design
mechanically detecting any nearby marijuana-policy change
(Figure~\ref{fig:coef}, bottom two rows).

\subsection{Persistence check}

Does the switcher effect continue, flatten, or reverse in the years
after a county crosses the 50-mile threshold? We re-estimate
Equation~\ref{eq:main} with event-time indicators for the switch
transition and each subsequent transition (one, and two-or-more,
transitions after switching), keeping switched counties in the panel
afterward purely for this check (they are otherwise dropped, per
Section~2). With the baseline-rate control, the switch-transition
effect remains significant ($-2.42$, $p=0.001$), but neither
subsequent-transition coefficient is significant (one transition
after: $+1.22$, $p=0.171$; two-plus: $+1.14$, $p=0.157$;
Figure~\ref{fig:persistence}), and a joint test rejects that the lag
effects equal the switch effect ($p=0.006$). The one-year decline
associated with switching does not continue to compound in later
years; if anything, the point estimates hint at partial reversion,
though not significantly so.

\begin{figure}[htbp]
  \centering
  \includegraphics[width=0.85\textwidth]{figures/fig6_persistence.png}
  \caption{Persistence check: switch-transition effect vs. subsequent-transition effects.}
  \label{fig:persistence}
\end{figure}

\subsection{Welfare translation}

National age-specific overdose rates (CDC NCHS Data Query System,
2024) imply a population-weighted average age at overdose death of
46.1 years. We translate the switcher coefficient into years of life
saved using the gap between life expectancy and this average age at
death, rather than the full remaining actuarial life expectancy, since
a person who dies of an overdose at 46 does not have the population's
full remaining life expectancy at stake, only the gap between when
they died and how long they could plausibly have lived absent this
specific cause. Life expectancy is drawn from CDC's state
life-expectancy-at-birth estimates \citep{CDCStateLE2021},
population-weighted across the 177 switcher counties' states (75.9
years), for a gap of 29.7 years per averted death. Applying the
preferred switcher coefficient to the 32.5 million residents of the
177 non-suppressed switcher counties implies approximately 755 averted
deaths (95\% CI: 245 to 1{,}265) over 2020--2025, and roughly 22,500
years of life saved (95\% CI: 7,300 to 37,600;
Figure~\ref{fig:agedist}). This is illustrative of what the point
estimate would mean, not a confirmed welfare effect: it inherits every
caveat on the underlying coefficient, including its greater
sensitivity to fixed-effects controls (Section~\ref{sec:pretrends}).

\begin{figure}[htbp]
  \centering
  \includegraphics[width=\textwidth]{figures/figC_age_distribution.png}
  \caption{National age distribution of drug-overdose deaths, 2024.}
  \label{fig:agedist}
\end{figure}

% ============================================================
% 4. CONCLUSION
% ============================================================
\section{Conclusion}

Extending the cross-border distance design to a stacked multi-period
panel that uses every state's actual legalization date resolves a
fragility in an earlier two-snapshot version: the continuous per-mile
specification is now sign-consistent with, and more robust than, the
binary switcher result, surviving even county fixed effects. A formal
pre-trends test finds and corrects a mean-reversion problem in the raw
switcher comparison; a persistence check finds the switch-year effect
does not compound in subsequent years; a medical-legalization placebo
check is null, as it should be if the effect is specific to
retail-market access; a comparison against canonical two-way
fixed-effects and Callaway--Sant'Anna estimators shows our stacked
design sitting between a TWFE estimate inflated in the direction its
known bias predicts and a Callaway--Sant'Anna estimate whose own
pre-treatment coefficients independently confirm the same
mean-reversion problem this paper corrects for -- adding our own
baseline-rate covariate directly to Callaway--Sant'Anna does not
resolve this and instead breaks down specifically on the single cohort
our own period-specific diagnostics already identify as imbalanced, at
every distance threshold we check, reinforcing rather than undermining
the case for the stacked design's explicit, cohort-aware correction; a
dispensary-level
robustness check that
replaces state-legalization timing with real, geocoded dispensary
opening dates for a 9-state subsample reproduces the same sign and a
comparable magnitude, though with weaker statistical power reflecting
its smaller, partial-coverage sample; and a check using already-near
counties as an alternative control group, on both treatment-timing
definitions, finds the switcher effect is never significant against
the state-legal-county measure and, against the dispensary measure,
significant without controls but not once mean reversion, state, or
county fixed effects are applied. Main
limitations: VSRR data are provisional and small-cell suppression
removes more than half the panel, concentrated where identifying
variation is richest; the 50-mile radius is a design choice; the panel
is bounded to 2020--2025; the welfare translation mixes a state-level
life-expectancy figure with a national age-at-death figure; the
dispensary-level check in Section~\ref{sec:dispensary} currently
covers only 9 of 24 legal states, so should be read as directionally
confirmatory rather than a fully independent replication; and the
switcher result's robustness to the choice of control group
(Section~\ref{sec:alwaysnear}) is weaker than its robustness to fixed
effects, so we lean more heavily on the continuous specification as
this paper's primary quantitative claim. We recommend this be treated
as suggestive evidence of an opioid-substitution channel specific to
recreational retail-market access, meriting confirmation with final,
non-provisional NCHS mortality data and, eventually, full
national dispensary-level coverage.

% ============================================================
% SUPPLEMENTARY MATERIAL
% ============================================================
\appendix
\section{Supplementary figures}
\label{sec:appendix}

\noindent\emph{Full robustness-check summary.} Table~\ref{tab:robustsummary}
consolidates every fixed-effects and sample-construction check reported
in this paper into one place. Panel A reproduces Table~\ref{tab:fe}
(the checks that bear directly on the pooled, headline specification
in Table~\ref{tab:main}). Panel B reports the additional checks used
only to diagnose and resolve the 2023--2024 period-specific sign
reversal discussed in Section~\ref{sec:periodeffects}; these are
estimated on that single transition in isolation ($n=691$) and were
never substituted into the pooled specification. The dispensary-level
robustness check (Section~\ref{sec:dispensary}, Table~\ref{tab:dispensary})
uses a different underlying sample (real dispensary-opening dates,
9-state coverage) and so is not reproduced in this table, which is
scoped to checks sharing the main $N=3{,}794$ sample.

\begin{table}[htbp]
\centering
\caption{Full robustness-check summary}
\label{tab:robustsummary}
\resizebox{\textwidth}{!}{%
\begin{tabular}{@{}lcccccc@{}}
\toprule
 & \multicolumn{3}{c}{Switcher (binary)} & \multicolumn{3}{c}{Distance (per mile)} \\
Specification & Coef. & $p$ (st.) & $p$ (cty.) & Coef. & $p$ (st.) & $p$ (cty.) \\
\midrule
\multicolumn{7}{@{}l}{\emph{Panel A: Pooled specification, all periods ($N=3{,}794$) -- bears on Table~\ref{tab:main}}} \\
No controls & $-3.151$ & 0.005 & $<0.001$ & $+0.0128^{***}$ & $<0.001$ & $<0.001$ \\
+ Baseline overdose-rate level (preferred) & $-2.320^{***}$ & 0.004 & 0.003 & $+0.0122^{***}$ & $<0.001$ & $<0.001$ \\
+ State fixed effects & $-1.646$ & 0.109 & 0.062 & $+0.0125^{***}$ & $<0.001$ & $<0.001$ \\
+ County fixed effects & $-0.999$ & 0.532 & 0.463 & $+0.0122^{**}$ & 0.023 & 0.009 \\
\addlinespace
\multicolumn{7}{@{}l}{\emph{Panel B: 2023--2024 transition only ($n=691$) -- diagnoses Section~\ref{sec:periodeffects}'s sign reversal}} \\
Baseline-rate control only & $+8.280^{***}$ & 0.008 & 0.022 & $+0.1132^{***}$ & $<0.001$ & $<0.001$ \\
+ Census-division fixed effects & $+2.786$ & 0.253 & 0.187 & $+0.0149$ & 0.335 & 0.294 \\
+ Winsorized baseline rate (95th pctile) & $+4.574^{**}$ & 0.032 & 0.011 & --- & --- & --- \\
+ Division FE \& winsorized baseline & $+2.328$ & 0.304 & 0.255 & $+0.0320^{**}$ & 0.036 & 0.044 \\
Excluding top 5\% by Cook's distance & $+1.822$ & 0.166 & 0.057 & --- & --- & --- \\
\bottomrule
\end{tabular}%
}
\vspace{2pt}
\begin{minipage}{0.95\textwidth}
\footnotesize\emph{Notes:} Panel A: each row a separate regression on
the identical $N=3{,}794$ pooled sample, 31 state / 1{,}249 county
clusters. Panel B: each row a separate regression restricted to the
2023--2024 transition only (691 rows); ``---'' denotes specifications
where the distance coefficient was not separately estimated. Because
this restricted subsample has one row per county (no repeat
observations within the transition), Panel B's ``county-clustered'' SE
is mechanically equivalent to a heteroskedasticity-robust SE with
singleton clusters, not a correction for within-cluster serial
correlation as in Panel A -- reported for completeness, not as a
like-for-like comparison to Panel A's county column. Panel B's raw,
unadjusted group means (no regression) show switcher counties declined
\emph{more} than controls in this transition ($-14.5$ vs.\ $-8.5$
deaths per 100,000); the positive sign in Panel B's first row appears
only once a linear baseline-rate control is applied to a period with
an unusually large baseline-rate imbalance between switchers and
controls (Section~\ref{sec:periodeffects}). Stars reflect the
state-clustered $p$-value. \emph{Correction:} Panel A's ``+ State fixed
effects'' row is corrected from $-1.411$/$+0.0129$ to $-1.646$/$+0.0125$
to match Table~\ref{tab:fe}; see that table's note. $^{**}p<0.05$,
$^{***}p<0.01$.
\end{minipage}
\end{table}

\noindent\emph{Always-near vs.\ always-far: a further control-group
check.} Table~\ref{tab:alwaysnear} above replaces the far/never-treated
control group with always-near counties while keeping switcher
counties as the treated group. Here we ask a different question:
setting switcher counties aside entirely, do counties that were
\emph{always} within 50 miles of a legal-market county throughout
2020--2025 (``always-near,'' 116 counties, never used as either
treatment or control elsewhere in this paper) look different from
counties that were \emph{always} 50 miles or farther away over the
same window (``always-far,'' 2,038 counties)? Neither group ever
switches, so we compare them using $\text{NearAccess}_i \in \{0,1\}$,
a time-invariant county characteristic, in the identical stacked
5-transition panel structure and escalating fixed-effects battery used
throughout the paper (transition FE, then the baseline-rate control,
then state FE, then census-division-by-transition and
state-by-transition FE) -- with one exception: we do not report a
county-fixed-effects row. Because $\text{NearAccess}_i$ never varies
within a county by construction, a county fixed effect would be
exactly collinear with it, leaving the coefficient structurally
unidentified rather than merely imprecise (verified directly: 0 of
1,153 sample counties have more than one $\text{NearAccess}$ value
across their rows, and forcing the regression anyway produces a design
matrix with condition number $\sim 3\times10^{17}$). This is a genuine
structural difference from the switcher-based checks elsewhere in the
paper, where the treatment variable changes within a county over the
panel and county fixed effects are identified.

\begin{table}[htbp]
\centering
\caption{Always-near vs.\ always-far counties (neither group ever switches)}
\label{tab:alwaysnearfar}
\begin{tabular}{@{}lcccc@{}}
\toprule
Specification & Coef. & $p$ (state) & $p$ (county) & $N$ \\
\midrule
No controls & $-0.431$ & 0.199 & 0.141 & 3{,}960 \\
+ Baseline overdose-rate level & $-0.596^{*}$ & 0.071 & 0.015 & 3{,}492 \\
+ State FE & $-0.367$ & 0.388 & 0.254 & 3{,}492 \\
+ Census division $\times$ transition FE & $+0.060$ & 0.847 & 0.837 & 3{,}492 \\
+ State $\times$ transition FE & $-0.118$ & 0.712 & 0.712 & 3{,}492 \\
\bottomrule
\end{tabular}
\vspace{2pt}
\begin{minipage}{0.92\textwidth}
\footnotesize\emph{Notes:} $\text{NearAccess}_i=1$ for the 116
always-near counties, $0$ for the 2,038 always-far counties (of which
890 have non-suppressed baseline-year data); switcher counties are
excluded entirely. 1,153 counties total (62 always-near, 1,091
always-far, contribute at least one non-suppressed row); 29 state
clusters throughout; county clusters $=$ 1,153 (No controls) or 890
(all other rows). No
county-fixed-effects row: see discussion above. $^{*}p<0.10$.
\end{minipage}
\end{table}

The coefficient is a consistently small, mostly null negative --
always-near counties show marginally lower overdose-mortality growth
than always-far counties once the baseline-rate control is added
($-0.60$, $p=0.071$ state-clustered, $p=0.015$ county-clustered), but
this does not survive any of the three fixed-effects specifications
that follow, all of which are far from significant under both
clustering choices. We read this as broadly consistent with the null
finding in Table~\ref{tab:alwaysnear}'s Panel A (switcher vs.\
always-near): whichever pair of the three groups (switcher, always-near,
always-far) is compared, the always-near counties do not show a
robust, controls-surviving difference from the rest of the sample --
a further reason we treat the continuous distance specification, not
any binary group comparison, as this paper's primary quantitative
claim.

\noindent\emph{Nonlinearity check.} Table~\ref{tab:main}'s continuous
specification imposes a constant marginal effect. We relax this by
adding quadratic and cubic terms in $\Delta\text{Dist}_{it}$ to the
preferred continuous specification (baseline overdose-rate control,
transition and county fixed effects, SE clustered by state, identical
$N=3{,}794$ sample). The quadratic term is significant on its own
($p=0.020$), the cubic term is marginal ($p=0.050$), and a joint
F-test that both nonlinear terms are zero rejects linearity
($F=5.41$, $p=0.010$). Figure~\ref{fig:doseresponse} plots the
resulting dose-response relationship (panel a: binned residualized
means with the fitted cubic overlaid) and the implied marginal effect
of a one-mile change in distance as a function of the dose itself
(panel b, with a 95\% confidence band from the delta method). The
marginal effect is negative and largest in magnitude for the biggest
distance reductions (around $-0.04$ deaths per 100,000 per mile at a
$-250$-mile dose) and flattens toward zero as the dose shrinks toward
no change -- consistent with the effect being driven by counties that
gain substantial proximity to a legal market, not by small
reshufflings of the nearest-market distance.

\begin{figure}[htbp]
  \centering
  \includegraphics[width=\textwidth]{figures/figA_doseresponse_multiperiod.png}
  \caption{Dose-response (panel a) and implied marginal effect (panel b), pooled stacked-transition panel.}
  \label{fig:doseresponse}
\end{figure}

\noindent\emph{Is there a natural discontinuity, and does it justify
50 miles specifically?} The binary switcher threshold used throughout
this paper (and its 70- and 100-mile alternatives, Appendix
Table~\ref{tab:thresh-main}) is a design choice, not something derived
from the mortality data itself. A natural check is the reverse
question: does overdose mortality, plotted against raw distance to the
nearest legal-market county, show a visible break \emph{anywhere}
between 50 and 100 miles that would justify treating counties on
either side of some line as ``treated'' and ``control,'' rather than
using the continuous measure throughout?

Figure~\ref{fig:distscan} answers this at coarse resolution: binned
mean overdose rates (10-mile bins) against raw distance, 0--250 miles,
for 2020 (thinnest legal-market coverage) and 2025 (widest), with the
50-, 70-, and 100-mile candidates marked. Neither year shows a sharp
break at any of the three candidates -- the relationship is a gradual
rise-then-plateau in 2020 and a gradual decline-then-flatten in 2025,
consistent in direction with this paper's continuous coefficient
(mortality lower nearer the border) but with no visually obvious kink
to hang a specific cutoff on.

\begin{figure}[htbp]
  \centering
  \includegraphics[width=\textwidth]{figures/figI_distance_scan.png}
  \caption{Binned mean overdose mortality vs.\ distance to the nearest legal-market county, 0--250 miles (10-mile bins), 2020 vs.\ 2025. Dashed lines mark the 50-, 70-, and 100-mile candidate thresholds used elsewhere in this paper. Hawaii excluded (distance not comparable). Marker size $\propto$ county count in the bin; solid line is a 3-bin rolling mean shown as a visual guide, not a fitted model.}
  \label{fig:distscan}
\end{figure}

A coarse scan can still miss a genuine \emph{local} jump, so
Figure~\ref{fig:rdplot50} zooms in with a formal CCT-style RD plot
\citep{CalonicoCattaneoTitiunik2014} centered exactly on the 50-mile
cutoff, restricted to $\pm$200 miles and fit with a local-linear
(not global-polynomial) trend on each side, same two years. Here
2020 \emph{does} show a visible jump -- the fitted level drops
sharply immediately to the left of the cutoff and steps back up just
past it -- matching the formal RD estimate at 50 miles for that
year ($+15.79$, $p=0.032$, CCT-optimal bandwidth, clustered by state).
2025 shows a much smaller, less discrete step, matching that year's
null RD estimate ($+1.10$, $p=0.82$). Differencing the two
(diff-in-disc, following \citet{GrembiNanniciniTroiano2016}) gives
$-14.68$ (se$=8.79$, $p=0.095$): the one clear break we can find is
specific to 2020 -- the baseline year, before almost any rollout had
occurred, so not itself interpretable as a treatment effect -- and it
shrinks rather than strengthens as the legal-market frontier expands
through 2025.

\begin{figure}[htbp]
  \centering
  \includegraphics[width=\textwidth]{figures/figJ_rdplot_50mi.png}
  \caption{Formal regression-discontinuity plot at the 50-mile cutoff, CCT-optimal bins and a local-linear fit on each side \citep{CalonicoCattaneoTitiunik2014}, restricted to $\pm$200 miles, 2020 vs.\ 2025. Marker size $\propto$ bin count. Vertical line is the 50-mile cutoff.}
  \label{fig:rdplot50}
\end{figure}

Taken together, these two figures are why this paper leans on the
continuous distance measure rather than asserting a specific
``true'' discontinuity point: the coarse scan finds no candidate
cutoff between 50 and 100 miles that stands out visually, and the one
place a formal local test does detect a break (50 miles, 2020) does not
survive as the legal-market frontier grows over the panel. We read
this as a reason for some caution about the near/far binary comparison
generally, not as confirming or overturning the main results -- and as
support for treating 50 miles as the pragmatic, sample-size-driven
choice described above (Section~\ref{sec:bandwidth-choice}), rather
than as a data-discovered break.

\noindent\emph{State-level trends and legalization timing.}
Figure~\ref{fig:statetrends} plots monthly overdose mortality,
2015--2026, for all 24 states in the source timing table, with both
the policy-effective date and the first-retail-sales date marked.
Figure~\ref{fig:eventstudy} pools the 23 states with a retail date,
aligned by months since first retail sales; neither figure feeds into
the county-level regression design.

\begin{figure}[htbp]
  \centering
  \includegraphics[width=\textwidth]{figures/figD_state_overdose_trends_new.png}
  \caption{State overdose mortality trends, 2015--2026, with policy-effective and retail-sales dates.}
  \label{fig:statetrends}
\end{figure}

\begin{figure}[htbp]
  \centering
  \includegraphics[width=\textwidth]{figures/figE_eventstudy_overdose_new.png}
  \caption{Overdose mortality around the retail-legalization cutoff, pooled across 23 states.}
  \label{fig:eventstudy}
\end{figure}

\noindent\emph{Robustness to the choice of distance threshold (70 and 100 miles).}
\label{sec:bandwidth-choice}
Every specification in this paper uses a 50-mile switcher/control
threshold, a design choice inherited from earlier work in this project
rather than a formally optimized cutoff (an earlier CCT MSE-optimal
bandwidth check landed close to 50 miles, but as confirmation of an
already-chosen threshold, not as the method that selected it). Here we
reproduce the main design (Table~\ref{tab:fe}), the dispensary-level
design (Table~\ref{tab:dispensary}), both panels of the
always-near control-group check (Table~\ref{tab:alwaysnear}), and the
TWFE/Callaway--Sant'Anna estimator comparison (Table~\ref{tab:twfecs})
at two alternative thresholds, 70 and 100 miles, using the identical
construction throughout -- only the cutoff changes.

\begin{table}[htbp]
\centering
\caption{Main design at alternative thresholds (cf.\ Table~\ref{tab:fe})}
\label{tab:thresh-main}
\resizebox{\textwidth}{!}{%
\begin{tabular}{@{}lcccc@{}}
\toprule
 & \multicolumn{2}{c}{70-mile threshold} & \multicolumn{2}{c}{100-mile threshold} \\
Specification & Switcher & Distance & Switcher & Distance \\
\midrule
No controls & \shortstack{$-3.843^{***}$\\{\scriptsize $p_{cty}<0.001$}} & \shortstack{$+0.0126^{***}$\\{\scriptsize $p_{cty}<0.001$}} & \shortstack{$-4.946^{**}$\\{\scriptsize $p_{cty}<0.001$}} & \shortstack{$+0.0129^{***}$\\{\scriptsize $p_{cty}<0.001$}} \\[1.2em]
+ Baseline overdose-rate level & \shortstack{$-3.022^{***}$\\{\scriptsize $p_{cty}<0.001$}} & \shortstack{$+0.0121^{***}$\\{\scriptsize $p_{cty}<0.001$}} & \shortstack{$-4.221^{**}$\\{\scriptsize $p_{cty}<0.001$}} & \shortstack{$+0.0127^{***}$\\{\scriptsize $p_{cty}<0.001$}} \\[1.2em]
+ State FE & \shortstack{$-2.402^{**}$\\{\scriptsize $p_{cty}=0.006$}} & \shortstack{$+0.0119^{***}$\\{\scriptsize $p_{cty}<0.001$}} & \shortstack{$-3.730^{*}$\\{\scriptsize $p_{cty}=0.002$}} & \shortstack{$+0.0122^{***}$\\{\scriptsize $p_{cty}<0.001$}} \\[1.2em]
+ Census division $\times$ transition FE & \shortstack{$-2.333^{**}$\\{\scriptsize $p_{cty}=0.010$}} & \shortstack{$+0.0084^{***}$\\{\scriptsize $p_{cty}=0.012$}} & \shortstack{$-3.685$\\{\scriptsize $p_{cty}=0.009$}} & \shortstack{$+0.0085^{***}$\\{\scriptsize $p_{cty}=0.012$}} \\[1.2em]
+ State $\times$ transition FE & \shortstack{$+0.677$\\{\scriptsize $p_{cty}=0.596$}} & \shortstack{$+0.0013$\\{\scriptsize $p_{cty}=0.775$}} & \shortstack{$-1.659^{*}$\\{\scriptsize $p_{cty}=0.190$}} & \shortstack{$+0.0005$\\{\scriptsize $p_{cty}=0.917$}} \\[1.2em]
+ County FE & \shortstack{$-1.888$\\{\scriptsize $p_{cty}=0.152$}} & \shortstack{$+0.0110^{**}$\\{\scriptsize $p_{cty}=0.024$}} & \shortstack{$-3.631$\\{\scriptsize $p_{cty}=0.038$}} & \shortstack{$+0.0109^{*}$\\{\scriptsize $p_{cty}=0.031$}} \\[1.2em]
\bottomrule
\end{tabular}%
}
\vspace{2pt}
\begin{minipage}{0.98\textwidth}
\footnotesize\emph{Notes:} $N=3{,}578$ at 70 miles, $N=3{,}291$ at 100
miles (231 and 298 raw switch events, respectively, vs.\ 177 at 50
miles). Stars reflect the state-clustered $p$-value; each cell's second
line reports the county-clustered $p$-value ($p_{cty}$) for the same
coefficient. $^{*}p<0.10$, $^{**}p<0.05$, $^{***}p<0.01$.
\end{minipage}
\end{table}

\begin{table}[htbp]
\centering
\caption{Dispensary-level design at alternative thresholds (cf.\ Table~\ref{tab:dispensary})}
\label{tab:thresh-dispensary}
\resizebox{\textwidth}{!}{%
\begin{tabular}{@{}lcccc@{}}
\toprule
 & \multicolumn{2}{c}{70-mile threshold} & \multicolumn{2}{c}{100-mile threshold} \\
Specification & Distance & Switcher & Distance & Switcher \\
\midrule
No controls & \shortstack{$+0.0230^{*}$\\{\scriptsize $p_{cty}<0.001$}} & \shortstack{$-4.648^{***}$\\{\scriptsize $p_{cty}<0.001$}} & \shortstack{$+0.0231^{*}$\\{\scriptsize $p_{cty}<0.001$}} & \shortstack{$-7.245^{***}$\\{\scriptsize $p_{cty}<0.001$}} \\[1.2em]
+ Baseline overdose-rate level & \shortstack{$+0.0172^{*}$\\{\scriptsize $p_{cty}<0.001$}} & \shortstack{$-2.457^{**}$\\{\scriptsize $p_{cty}=0.041$}} & \shortstack{$+0.0168^{*}$\\{\scriptsize $p_{cty}<0.001$}} & \shortstack{$-4.776^{**}$\\{\scriptsize $p_{cty}<0.001$}} \\[1.2em]
+ State FE & \shortstack{$+0.0180$\\{\scriptsize $p_{cty}<0.001$}} & \shortstack{$-1.828$\\{\scriptsize $p_{cty}=0.166$}} & \shortstack{$+0.0174$\\{\scriptsize $p_{cty}<0.001$}} & \shortstack{$-4.502^{*}$\\{\scriptsize $p_{cty}=0.004$}} \\[1.2em]
+ Census division $\times$ transition FE & \shortstack{$+0.0169$\\{\scriptsize $p_{cty}<0.001$}} & \shortstack{$-1.615$\\{\scriptsize $p_{cty}=0.323$}} & \shortstack{$+0.0165$\\{\scriptsize $p_{cty}<0.001$}} & \shortstack{$-4.287^{**}$\\{\scriptsize $p_{cty}=0.032$}} \\[1.2em]
+ State $\times$ transition FE & \shortstack{$+0.0046$\\{\scriptsize $p_{cty}=0.233$}} & \shortstack{$+4.569$\\{\scriptsize $p_{cty}=0.057$}} & \shortstack{$+0.0044$\\{\scriptsize $p_{cty}=0.249$}} & \shortstack{$+3.030$\\{\scriptsize $p_{cty}=0.219$}} \\[1.2em]
+ County FE & \shortstack{$+0.0212$\\{\scriptsize $p_{cty}<0.001$}} & \shortstack{$-1.603$\\{\scriptsize $p_{cty}=0.352$}} & \shortstack{$+0.0204$\\{\scriptsize $p_{cty}<0.001$}} & \shortstack{$-5.419^{*}$\\{\scriptsize $p_{cty}=0.010$}} \\[1.2em]
\bottomrule
\end{tabular}%
}
\vspace{2pt}
\begin{minipage}{0.98\textwidth}
\footnotesize\emph{Notes:} $N=3{,}852$ at 70 miles, $N=3{,}678$ at 100
miles (124 and 151 raw switch events, respectively, vs.\ 91 at 50
miles). Stars reflect the state-clustered $p$-value; each cell's second
line reports the county-clustered $p$-value ($p_{cty}$). $^{*}p<0.10$,
$^{**}p<0.05$, $^{***}p<0.01$.
\end{minipage}
\end{table}

\begin{table}[htbp]
\centering
\caption{Always-near control-group check at alternative thresholds (cf.\ Table~\ref{tab:alwaysnear})}
\label{tab:thresh-alwaysnear}
\resizebox{\textwidth}{!}{%
\begin{tabular}{@{}lcccc@{}}
\toprule
 & \multicolumn{2}{c}{70-mile threshold} & \multicolumn{2}{c}{100-mile threshold} \\
Specification & Panel A (state) & Panel B (dispensary) & Panel A (state) & Panel B (dispensary) \\
\midrule
No controls & \shortstack{$-2.542^{**}$\\{\scriptsize $p_{cty}=0.003$}} & \shortstack{$-1.960$\\{\scriptsize $p_{cty}=0.298$}} & \shortstack{$-3.944^{**}$\\{\scriptsize $p_{cty}<0.001$}} & \shortstack{$-5.505^{***}$\\{\scriptsize $p_{cty}<0.001$}} \\[1.2em]
+ Baseline overdose rate & \shortstack{$-2.819^{***}$\\{\scriptsize $p_{cty}<0.001$}} & \shortstack{$-0.286$\\{\scriptsize $p_{cty}=0.882$}} & \shortstack{$-3.838^{**}$\\{\scriptsize $p_{cty}<0.001$}} & \shortstack{$-3.567^{*}$\\{\scriptsize $p_{cty}=0.025$}} \\[1.2em]
+ State FE & \shortstack{$-2.659^{**}$\\{\scriptsize $p_{cty}=0.009$}} & \shortstack{$-0.072$\\{\scriptsize $p_{cty}=0.973$}} & \shortstack{$-3.801^{*}$\\{\scriptsize $p_{cty}=0.003$}} & \shortstack{$-3.932^{*}$\\{\scriptsize $p_{cty}=0.023$}} \\[1.2em]
+ Census division $\times$ transition FE & \shortstack{$-2.134$\\{\scriptsize $p_{cty}=0.043$}} & \shortstack{$+0.715$\\{\scriptsize $p_{cty}=0.700$}} & \shortstack{$-2.470^{*}$\\{\scriptsize $p_{cty}=0.031$}} & \shortstack{$+0.664$\\{\scriptsize $p_{cty}=0.747$}} \\[1.2em]
+ State $\times$ transition FE & \shortstack{$-1.862$\\{\scriptsize $p_{cty}=0.179$}} & \shortstack{$+2.986$\\{\scriptsize $p_{cty}=0.200$}} & \shortstack{$-1.648$\\{\scriptsize $p_{cty}=0.212$}} & \shortstack{$+2.871$\\{\scriptsize $p_{cty}=0.313$}} \\[1.2em]
+ County FE & \shortstack{$-2.612$\\{\scriptsize $p_{cty}=0.075$}} & \shortstack{$+0.940$\\{\scriptsize $p_{cty}=0.743$}} & \shortstack{$-3.869$\\{\scriptsize $p_{cty}=0.033$}} & \shortstack{$-4.831^{*}$\\{\scriptsize $p_{cty}=0.049$}} \\[1.2em]
\bottomrule
\end{tabular}%
}
\vspace{2pt}
\begin{minipage}{0.98\textwidth}
\footnotesize\emph{Notes:} Binary switcher specification only, both
panels, at both thresholds, matching Table~\ref{tab:alwaysnear}'s
scope. At 70 miles: Panel A $N=1{,}033$ (no controls) / $961$ (all
other rows), 231 ever-switcher + 98 always-near counties; Panel B
$N=702$/$666$, 124 ever-switcher + 59 always-near counties. At 100
miles: Panel A $N=1{,}377$/$1{,}260$, 298 ever-switcher + 151
always-near counties; Panel B $N=979$/$931$, 151 ever-switcher + 108
always-near counties. Stars reflect the state-clustered $p$-value;
each cell's second line reports the county-clustered $p$-value
($p_{cty}$). $^{*}p<0.10$, $^{**}p<0.05$, $^{***}p<0.01$.
\end{minipage}
\end{table}

The main design (Table~\ref{tab:thresh-main}) is strikingly stable
across all three thresholds: the continuous coefficient stays in the
$+0.008$ to $+0.013$ range throughout, significant in every
specification except the same sparse state-by-transition row that
also fails at 50 miles, and the switcher coefficient shows the same
qualitative pattern of weakening under state and county fixed effects
at every threshold. If anything, both measures are slightly
\emph{more} significant at 100 miles than at 50, since the larger
radius produces more switch events (298 vs.\ 177) without changing the
underlying relationship. The dispensary-level design
(Table~\ref{tab:thresh-dispensary}) is similarly stable in sign and
magnitude, though the switcher coefficient actually strengthens at 100
miles (surviving state, division, and county fixed effects, unlike at
50 or 70 miles) -- consistent with the dispensary check's smaller
sample simply gaining power as the radius widens, rather than the
relationship itself changing. The always-near control-group check
(Table~\ref{tab:thresh-alwaysnear}) shows the same qualitative fragility
at all three thresholds: Panel A (state-legal-county distance) never
reaches significance under the stricter fixed-effects rows at any
threshold, and Panel B (dispensary distance) is sensitive to controls
throughout, including one sign flip under census division-by-transition
fixed effects at 100 miles that mirrors the same low-power pattern
already flagged for the analogous 50-mile row. Overall, we read this as
evidence that 50 miles was a reasonable choice: the paper's qualitative
conclusions -- the continuous specification is more robust than the
switcher indicator, and the switcher result is more fragile to the
choice of control group than to fixed effects -- hold at 70 and 100
miles as well, rather than being an artifact of the specific radius
chosen.

\begin{table}[htbp]
\centering
\caption{TWFE/Callaway--Sant'Anna estimator comparison at alternative thresholds (cf.\ Table~\ref{tab:twfecs})}
\label{tab:thresh-twfecs}
\resizebox{\textwidth}{!}{%
\footnotesize
\begin{tabular}{@{}lcccccccccc@{}}
\toprule
 & \multicolumn{5}{c}{70-mile threshold} & \multicolumn{5}{c}{100-mile threshold} \\
 & Coef. & \multicolumn{2}{c}{State-clustered} & \multicolumn{2}{c}{County-clustered} & Coef. & \multicolumn{2}{c}{State-clustered} & \multicolumn{2}{c}{County-clustered} \\
Estimator & (70mi) & SE & $p$ & SE & $p$ & (100mi) & SE & $p$ & SE & $p$ \\
\midrule
TWFE (county + year FE) & $-6.400$ & $2.173$ & $0.003$ & $1.124$ & $<0.001$ & $-7.293$ & $2.878$ & $0.011$ & $1.197$ & $<0.001$ \\
Callaway--Sant'Anna (CS) & $-2.843$ & $1.856$ & $0.126$ & $1.011$ & $0.005$ & $-4.078$ & $2.487$ & $0.101$ & $1.126$ & $<0.001$ \\
\quad + baseline-rate covariate & $+10.134$ & $6.159$ & $0.100$ & $4.468$ & $0.023$ & $-0.408$\rlap{$^{\dagger}$} & $0.922$ & $0.658$ & $0.653$ & $0.532$ \\
\bottomrule
\end{tabular}%
}
\vspace{2pt}
\begin{minipage}{0.95\textwidth}
\footnotesize\emph{Notes:} CS uses never-treated controls throughout.
TWFE and CS use the balanced panel at each threshold ($n=3{,}570$ at 70
miles, 594 county clusters; $n=3{,}432$ at 100 miles, 572 county
clusters). Column headers are laid out as Coef./SE/$p$ under each
clustering level, matching Table~\ref{tab:twfecs}.
$^{\dagger}$At 100 miles the 2024 cohort's group-time cells are not
estimable under the covariate specification (a doubly-robust
estimation failure, not a coding error) and are dropped from the
aggregate; see discussion in Section~\ref{sec:twfecs}.
\end{minipage}
\end{table}

The covariate-adjusted Callaway--Sant'Anna instability documented in
Section~\ref{sec:twfecs} is not specific to 50 miles: at 70 miles the
covariate-adjusted overall ATT again flips sign and stays only
marginally short of significance ($+10.13$, $p=0.10$); at 100 miles
the 2024 cohort's imbalance is severe enough that the estimator cannot
compute its group-time cells at all, and dropping that cohort leaves a
small, correctly-signed, imprecise estimate ($-0.41$, $p=0.66$). The
uncorrected TWFE and Callaway--Sant'Anna rows, by contrast, are stable
in sign and pattern across all three thresholds -- TWFE roughly double
the stacked design's magnitude throughout, plain Callaway--Sant'Anna
consistently smaller and insignificant. Only the covariate-adjusted row
is threshold-sensitive, and always in the same way: it breaks down on
whichever cohort has the worst baseline-rate overlap with its controls,
which is the 2024/Ohio-border cohort at every threshold we check.

% ============================================================
% REFERENCES
% ============================================================
\bibliographystyle{apalike}
\bibliography{references}

\end{document}
